\section{Definitions and invariance}

Let $X$ be a set of size $n$.  A \emph{Latin square} on $X$ is a map
$L\colon X\times X\to X$ such that every row and every column is a bijection.
For $a,c,u\in X$, define the row, column, and symbol-line bijections
\[
 r_a(c)=L(a,c),\qquad
 q_c(a)=L(a,c),\qquad
 s_u(c)=\text{the unique }a\text{ such that }L(a,c)=u.
\]
For two lines in the corresponding view, put
\[
 \rho_{a,b}=r_b^{-1}r_a,\qquad
 \kappa_{c,d}=q_d^{-1}q_c,\qquad
 \sigma_{u,v}=s_v^{-1}s_u.
\]
All three permutations act on an $n$-element complementary coordinate set.

\begin{lemma}\label{lem:fixed-point-free}
If the two lines are distinct, the corresponding induced permutation is
fixed-point-free.
\end{lemma}

\begin{proof}
A fixed point of $\rho_{a,b}$ would give $L(a,c)=L(b,c)$ in one column,
contrary to the Latin property.  The column case is identical.  A fixed point
of $\sigma_{u,v}$ would say that the same cell contains the two distinct
symbols $u$ and $v$.
\end{proof}

\begin{definition}
A view of $L$ \emph{fails} if one of its induced two-line permutations has an
odd cycle of length greater than one.  Write
\[
 \pat(L)=(\rowv(L),\colv(L),\symv(L))\in\{F,T\}^3,
\]
where $T$ means that the view fails.  A square is \emph{FFF} if
$\pat(L)=\FFF$.
\end{definition}

The terminology ``fails'' is chosen from the perspective of the local-trade
program: a $T$ view supplies an odd-cycle trade, while an $F$ view does not.

\begin{definition}
Fix two rows $a\ne b$ and a union $C$ of cycles of $\rho_{a,b}$.  The
associated \emph{two-line trade} swaps the entries in rows $a,b$ in precisely
the columns in $C$.  Column- and symbol-view trades are defined by permuting
the coordinate roles.
\end{definition}

\begin{proposition}\label{prop:trade}
A two-line trade produces another Latin square.  Define the three line-sign
products
\[
 R(L)=\prod_a\operatorname{sgn}(r_a),\qquad
 C(L)=\prod_c\operatorname{sgn}(q_c),\qquad
 S(L)=\prod_u\operatorname{sgn}(s_u),
\]
and put
\[
 \epsilon_{RC}=RC,\qquad \epsilon_{RS}=RS,\qquad \epsilon_{CS}=CS.
\]
The first is the conventional row--column Latin sign, while the other two
are its parastrophic companions.  If a trade has support size $k$ and
$\delta=(-1)^k$, then its sign ratios are
\[
\begin{array}{c|ccc|ccc}
\text{trade view} & R'/R&C'/C&S'/S&
 \epsilon'_{RC}/\epsilon_{RC}&\epsilon'_{RS}/\epsilon_{RS}&
 \epsilon'_{CS}/\epsilon_{CS}\\
\hline
\rowv    &1&\delta&\delta&\delta&\delta&1\\
\colv    &\delta&1&\delta&\delta&1&\delta\\
\symv    &\delta&\delta&1&1&\delta&\delta
\end{array}
\]
Consequently an odd row- or column-view trade reverses the conventional
$\epsilon_{RC}$, whereas an odd symbol-view trade preserves $\epsilon_{RC}$
and reverses $\epsilon_{RS}$ and $\epsilon_{CS}$.
\end{proposition}

\begin{proof}
Write $\pi=\rho_{a,b}$.  Since $r_a(c)=r_b(\pi(c))$ and $C$ is
$\pi$-invariant, the two rows contain the same symbol set on $C$.  Swapping
those entries therefore preserves both row symbol sets, while each affected
column merely swaps two entries.  Thus the result is Latin.

Extend $\pi|_C$ by the identity outside $C$.  The two changed row
permutations are obtained by precomposition with $\pi$ and $\pi^{-1}$, so
their sign changes cancel and $R'/R=1$.  Each of the $|C|$ affected column
permutations is composed with a transposition, so $C'/C=(-1)^{|C|}$.  The
same support contains the same $|C|$ symbols in the two rows; each associated
symbol-position permutation swaps those two rows, giving
$S'/S=(-1)^{|C|}$.  This proves the first table row.  Exchanging rows and
columns proves the second.

For a symbol-view trade, each support column swaps the two selected symbols,
and the cycle relation pairs the affected cells in exactly the same number of
rows.  Hence $R'/R=C'/C=(-1)^{|C|}$.  The two changed symbol-position
permutations differ from their predecessors by mutually inverse
permutations on the support, so their sign changes cancel and $S'/S=1$.
Multiplying the appropriate line-sign ratios proves the last three columns.
\end{proof}

\begin{proposition}[Tripartite link formulation]
\label{prop:tripartite-links}
Regard the cells of $L$ as the triangle decomposition
\[
 \bigl\{\{a,c,L(a,c)\}:a,c\in X\bigr\}
\]
of the complete tripartite graph on the row, column, and symbol copies of
$X$.  The links of the vertices in each part form a 1-factorization between
the other two parts.  An induced two-line cycle of length $\ell$ corresponds
bijectively to an alternating cycle of length $2\ell$ in the union of the two
associated link factors.  Hence $L$ is FFF if and only if every two-factor
union in all three link factorizations has only cycle lengths divisible by
four.

Furthermore, the total number of $4$-cycles among the two-factor unions is
the same in all three views and equals the number of intercalates of $L$.
\end{proposition}

\begin{proof}
For a row $a$, its incident triangles give the perfect matching
$M_a=\{\{c,r_a(c)\}:c\in X\}$ between columns and symbols.  For distinct
rows $a,b$, alternately following an edge of $M_a$ and an edge of $M_b$
sends a column $c$ to $r_b^{-1}r_a(c)=\rho_{a,b}(c)$.  Thus an $\ell$-cycle
of $\rho_{a,b}$ is exactly one alternating $2\ell$-cycle of
$M_a\cup M_b$.  Cyclically permuting the three coordinate classes proves the
column and symbol statements and the FFF equivalence.

A $4$-cycle in $M_a\cup M_b$ is exactly a $2$-cycle of $\rho_{a,b}$, hence a
$2$-by-$2$ Latin subsquare on two rows, two columns, and two symbols.  This is
an intercalate.  The same intercalate is counted once by its row pair, once by
its column pair, and once by its symbol pair, proving the three equalities.
\end{proof}

\begin{lemma}[Global three-view sign identity]
\label{lem:global-three-view-sign}
For every Latin square of order $n$,
\[
 R(L)C(L)S(L)=(-1)^{\binom n2}.
\]
\end{lemma}

\begin{proof}
Give each coordinate copy of $X$ the same fixed order, and let
$\tau(x,y)=(y,x)$ on $X^2$.  The $\binom n2$ off-diagonal pairs are the
transpositions of $\tau$, so
$\operatorname{sgn}(\tau)=(-1)^{\binom n2}$.

The row-entry map $A(r,c)=(r,L(r,c))$ is block diagonal by rows and has sign
$R(L)$.  The map $B(r,c)=(c,L(r,c))$ is a coordinate transposition followed
by the block-diagonal family of column permutations, so its sign is
$(-1)^{\binom n2}C(L)$.  Likewise, $D(c,s)=(s,s_s(c))$ has sign
$(-1)^{\binom n2}S(L)$.  Since $A=\tau\circ D\circ B$, taking signs gives
\[
 R(L)=(-1)^{3\binom n2}S(L)C(L)
      =(-1)^{\binom n2}S(L)C(L),
\]
which is equivalent to the claimed identity.
\end{proof}

For later comparison, let $C_{v,k}$ denote the total number of $k$-cycles
among all induced line-pair permutations in view $v$.

\begin{proposition}[Aggregate cycle-profile baseline]
\label{prop:aggregate-profile-baseline}
Let $n$ be even and let $p$ be prime.  If
\[
 \sum_{k=2}^n a_kx_k+b=0\pmod p
\]
for every cycle-count vector $(x_2,\ldots,x_n)$ of a fixed-point-free
permutation of degree $n$, then, separately for each view $v$,
\[
 \sum_{k=2}^n a_k C_{v,k}+\binom n2 b=0\pmod p.
\]
Moreover,
\[
 C_{\rowv,2}=C_{\colv,2}=C_{\symv,2}
\]
and
\[
 \sum_{v\in\{\rowv,\colv,\symv\}}\ \sum_{\substack{2\leq k\leq n\\k\text{ even}}}
 C_{v,k}=\frac{n(n-1)}2\pmod2.
\]
\end{proposition}

\begin{proof}
The aggregate profile in one view is the sum of the
$\binom n2$ fixed-point-free line-pair profiles, proving the first identity.
The second identity is the intercalate bijection from
Proposition~\ref{prop:tripartite-links}.  For the last identity, the sign of
a fixed-point-free permutation of even degree is $(-1)^e$, where $e$ is its
number of even cycles.  Multiplying relative-permutation signs over all line
pairs in a view gives the product of its line signs.
Lemma~\ref{lem:global-three-view-sign} gives the product over the three views as
$(-1)^{n(n-1)/2}$, which proves the congruence.
\end{proof}

\begin{proposition}[Line-pair sign cut]\label{prop:line-sign-cut}
Fix one view and write its line bijections as $\lambda_i$.  For
$P_{ij}=\lambda_j^{-1}\lambda_i$, the pairs for which $P_{ij}$ is odd form
a cut of the complete graph on the lines.  If $n$ is even and the view is
F, then the pairs for which $P_{ij}$ has an odd number of cycles form the
same cut.  Consequently, for even $n>2$, no view can have every induced
line-pair permutation equal to one $n$-cycle.
\end{proposition}

\begin{proof}
Put $\varepsilon_i=\operatorname{sgn}(\lambda_i)$.  Since sign is a
homomorphism,
\[
 \operatorname{sgn}(P_{ij})=\varepsilon_i\varepsilon_j.
\]
Thus $P_{ij}$ is odd exactly when the endpoint signs differ.  These pairs
are the edges of the cut between positive and negative lines.

In an F view of even degree, every cycle of $P_{ij}$ has even length and
therefore negative sign.  Hence
$\operatorname{sgn}(P_{ij})=(-1)^{c(P_{ij})}$, where $c(P_{ij})$ is the
number of cycles.  Finally, if all relative permutations were $n$-cycles,
every edge would be negative.  A triangle of three lines would then cross
the sign cut three times, whereas every cut meets a cycle evenly.
\end{proof}

An \emph{isotopy} independently permutes rows, columns, and symbols.  A
\emph{parastrophy} permutes the three coordinate roles.  Their combined
orbits are the \emph{main classes} (also called species).

\begin{proposition}\label{prop:invariance}
Under isotopy, each of the three view bits is invariant.  Under parastrophy,
the three bits are permuted.  Consequently, being FFF and having at least one
failing view are main-class invariants.
\end{proposition}

\begin{proof}
Let an isotopy act by permutations $\alpha,\beta,\gamma$ of rows, columns,
and symbols.  The transformed row maps satisfy
$r'_{\alpha(a)}=\gamma r_a\beta^{-1}$, and hence
\[
 (r'_{\alpha(b)})^{-1}r'_{\alpha(a)}
   =\beta(r_b^{-1}r_a)\beta^{-1}.
\]
Thus row cycle lengths are unchanged.  The column calculation conjugates by
$\alpha$, and the symbol calculation conjugates by $\beta$.  Parastrophy
only permutes which coordinate pair is called row, column, or symbol.
\end{proof}
